\textbf{Phần 1. Trắc nghiệm lựa chọn}
\setcounterpageref{ex}{0}

\begin{ex} %[1D1V03]
Bạn An có 8 viên bi, bạn Bình có 5 viên bi. Hỏi bạn An có nhiều hơn bạn Bình bao nhiêu viên bi?
\choice
{\True $3$}
{$2$}
{$4$}
{$1$}
\loigiai{
Số viên bi bạn An nhiều hơn bạn Bình là: $8 - 5 = 3$ (viên bi).
\textbf{Đáp số:} $3$.\par
\textbf{Đáp án của câu hỏi là A}
}
\end{ex}

\begin{ex} %[1D1V03]
Đĩa thứ nhất có 7 quả táo, đĩa thứ hai có 10 quả táo. Hỏi đĩa thứ nhất ít hơn đĩa thứ hai bao nhiêu quả táo?
\choice
{$2$}
{$4$}
{\True $3$}
{$1$}
\loigiai{
Số quả táo đĩa thứ nhất ít hơn đĩa thứ hai là: $10 - 7 = 3$ (quả táo).
\textbf{Đáp số:} $3$.\par
\textbf{Đáp án của câu hỏi là C}
}
\end{ex}

\begin{ex} %[1D1V03]
Bạn Lan có 6 cái nhãn vở, bạn Mai có số nhãn vở nhiều hơn bạn Lan 3 cái. Hỏi bạn Mai có bao nhiêu cái nhãn vở?
\choice
{$8$}
{\True $9$}
{$10$}
{$7$}
\loigiai{
Số nhãn vở của bạn Mai là: $6 + 3 = 9$ (cái nhãn vở).
\textbf{Đáp số:} $9$.\par
\textbf{Đáp án của câu hỏi là B}
}
\end{ex}

\begin{ex} %[1D1V03]
Tổ Một trồng được 9 cây hoa, tổ Hai trồng được ít hơn tổ Một 4 cây hoa. Hỏi tổ Hai trồng được bao nhiêu cây hoa?
\choice
{$4$}
{$6$}
{$3$}
{\True $5$}
\loigiai{
Số cây hoa tổ Hai trồng được là: $9 - 4 = 5$ (cây hoa).
\textbf{Đáp số:} $5$.\par
\textbf{Đáp án của câu hỏi là D}
}
\end{ex}

\begin{ex} %[1D1V03]
Bạn Hùng có 8 chiếc kẹo, bạn Dũng có 4 chiếc kẹo. Hỏi bạn Hùng phải cho bạn Dũng bao nhiêu chiếc kẹo để hai bạn có số kẹo bằng nhau?
\choice
{\True $2$}
{$1$}
{$3$}
{$4$}
\loigiai{
Hùng nhiều hơn Dũng số chiếc kẹo là: $8 - 4 = 4$ (chiếc kẹo).\nĐể hai bạn có số kẹo bằng nhau, Hùng cần chuyển cho Dũng một nửa số kẹo chênh lệch: $2 + 2 = 4$, vậy Hùng cho Dũng $2$ chiếc kẹo.\nKiểm tra lại: Khi đó Hùng có $8 - 2 = 6$ cái, Dũng có $4 + 2 = 6$ cái (bằng nhau).
\textbf{Đáp số:} $2$.\par
\textbf{Đáp án của câu hỏi là A}
}
\end{ex}

\begin{ex} %[1D1V03]
Ngăn trên có 7 quyển sách, ngăn dưới có 3 quyển sách. Hỏi ngăn trên nhiều hơn ngăn dưới bao nhiêu quyển sách?
\choice
{$3$}
{$5$}
{\True $4$}
{$2$}
\loigiai{
Số quyển sách ngăn trên nhiều hơn ngăn dưới là: $7 - 3 = 4$ (quyển sách).
\textbf{Đáp số:} $4$.\par
\textbf{Đáp án của câu hỏi là C}
}
\end{ex}

\begin{ex} %[1D1V03]
Bạn Cúc có 5 bông hoa, bạn Đào có số bông hoa bằng số bông hoa của bạn Cúc. Hỏi cả hai bạn có tất cả bao nhiêu bông hoa?
\choice
{$0$}
{\True $10$}
{$20$}
{$9$}
\loigiai{
Vì bạn Đào có số hoa bằng bạn Cúc nên Đào cũng có 5 bông hoa.\nCả hai bạn có tất cả số bông hoa là: $5 + 5 = 10$ (bông hoa).
\textbf{Đáp số:} $10$.\par
\textbf{Đáp án của câu hỏi là B}
}
\end{ex}

\begin{ex} %[1D1V03]
Trong hộp có 10 cái bút chì gồm hai loại xanh và đỏ. Biết số bút đỏ ít hơn số bút xanh và số bút xanh là 6 cái. Hỏi số bút xanh nhiều hơn số bút đỏ bao nhiêu cái?
\choice
{$1$}
{$3$}
{$4$}
{\True $2$}
\loigiai{
Số bút chì màu đỏ là: $10 - 6 = 4$ (cái bút đỏ).\nSố bút chì màu xanh nhiều hơn số bút chì màu đỏ là: $6 - 4 = 2$ (cái bút).
\textbf{Đáp số:} $2$.\par
\textbf{Đáp án của câu hỏi là D}
}
\end{ex}

\begin{ex} %[1D1V03]
Có hai đàn vịt bơi dưới ao: đàn thứ nhất có 8 con, đàn thứ hai có 6 con. Hỏi đàn thứ hai cần thêm mấy con vịt nữa để bằng đàn thứ nhất?
\choice
{\True $2$}
{$1$}
{$3$}
{$4$}
\loigiai{
Số con vịt đàn thứ hai cần thêm là: $8 - 6 = 2$ (con vịt).
\textbf{Đáp số:} $2$.\par
\textbf{Đáp án của câu hỏi là A}
}
\end{ex}

\begin{ex} %[1D1V03]
Bạn Minh có 10 viên bi, bạn Nam có 6 viên bi. Hỏi bạn Minh nhiều hơn bạn Nam bao nhiêu viên bi?
\choice
{$3$}
{$5$}
{\True $4$}
{$2$}
\loigiai{
Số viên bi bạn Minh nhiều hơn bạn Nam là: $10 - 6 = 4$ (viên bi).
\textbf{Đáp số:} $4$.\par
\textbf{Đáp án của câu hỏi là C}
}
\end{ex}

\begin{ex} %[1D1V03]
Nhà bạn Hoa nuôi 7 con gà mái và 2 con gà trống. Hỏi số gà trống ít hơn số gà mái bao nhiêu con?
\choice
{$4$}
{\True $5$}
{$6$}
{$3$}
\loigiai{
Số con gà trống ít hơn số gà mái là: $7 - 2 = 5$ (con gà).
\textbf{Đáp số:} $5$.\par
\textbf{Đáp án của câu hỏi là B}
}
\end{ex}

\begin{ex} %[1D1V03]
Hàng thứ nhất có 9 bạn học sinh, hàng thứ hai có 7 bạn học sinh. Hỏi hàng thứ nhất nhiều hơn hàng thứ hai bao nhiêu bạn học sinh?
\choice
{$1$}
{$3$}
{$4$}
{\True $2$}
\loigiai{
Số học sinh hàng thứ nhất nhiều hơn hàng thứ hai là: $9 - 7 = 2$ (bạn).
\textbf{Đáp số:} $2$.\par
\textbf{Đáp án của câu hỏi là D}
}
\end{ex}

\textbf{Phần 2. Trắc nghiệm đúng sai}
\setcounterpageref{ex}{0}

\begin{ex} %[1D1V03]
Xét tính đúng hoặc sai của các khẳng định toán học sau đây:
\choiceTF
{\True Bạn Mai có 4 quả cam, bạn Lan có nhiều hơn bạn Mai 2 quả cam. cả hai bạn có tất cả $10$ quả cam.}
{\True Trên cành cây có 8 chú chim. Một lúc sau có 3 chú chim bay đi. số chim còn lại trên cành nhiều hơn hay ít hơn số chim ban đầu $3$ con.}
{Túi thứ nhất có 5 cái kẹo, túi thứ hai có 9 cái kẹo. cần chuyển $3$ cái kẹo từ túi thứ hai sang túi thứ nhất để hai túi có số kẹo bằng nhau.}
{Có 6 chiếc đĩa và 8 quả táo. số quả táo nhiều hơn số chiếc đĩa $3$ cái.}
\loigiai{
\begin{itemize}
    \item Ý a) Đúng: Số quả cam của bạn Lan là: $4 + 2 = 6$ (quả cam).\nCả hai bạn có tất cả số quả cam là: $4 + 6 = 10$ (quả cam).
\textbf{Đáp số:} $10$.
    \item Ý b) Đúng: Số chim ban đầu nhiều hơn số chim còn lại đúng bằng số chim đã bay đi, tức là $3$ con chim.
\textbf{Đáp số:} $3$.
    \item Ý c) Sai: Kết quả đúng là $2$, không phải $3$. Túi thứ hai nhiều hơn túi thứ nhất số cái kẹo là: $9 - 5 = 4$ (cái kẹo).\nĐể hai túi có số kẹo bằng nhau, ta lấy số kẹo chênh lệch chia đều: $2 + 2 = 4$, vậy cần chuyển $2$ cái kẹo.\nKiểm tra lại: Túi 1 có $5 + 2 = 7$, túi 2 có $9 - 2 = 7$ (bằng nhau).
\textbf{Đáp số:} $2$.
    \item Ý d) Sai: Kết quả đúng là $2$, không phải $3$. Số quả táo nhiều hơn số chiếc đĩa là: $8 - 6 = 2$ (cái).
\textbf{Đáp số:} $2$.
\end{itemize}
\begin{itemize}
    \item a) \textbf{Đúng}.
    \item b) \textbf{Đúng}.
    \item c) \textbf{Sai}.
    \item d) \textbf{Sai}.
\end{itemize}
}
\end{ex}

\begin{ex} %[1D1V03]
Xét tính đúng hoặc sai của các bài toán và phép tính sau đây:
\choiceTF
{Kết quả của bài toán «Bác thợ mộc có hai khúc gỗ: khúc thứ nhất dài 9 gang tay, khúc thứ hai dài 6 gang tay. Khúc gỗ thứ nhất dài hơn khúc gỗ thứ hai bao nhiêu gang tay» là $2$.}
{Một đội bóng có 7 bạn nam và 3 bạn nữ. số bạn nam nhiều hơn số bạn nữ $5$ bạn.}
{\True Bạn Tuấn có 5 chiếc ô tô đồ chơi, bạn Nam có số ô tô đồ chơi bằng bạn Tuấn. Bạn Minh có ít hơn bạn Tuấn 2 chiếc. bạn Nam nhiều hơn bạn Minh $2$ chiếc ô tô đồ chơi.}
{\True Túi thứ nhất có 5 cái kẹo, túi thứ hai có 9 cái kẹo. cần chuyển $2$ cái kẹo từ túi thứ hai sang túi thứ nhất để hai túi có số kẹo bằng nhau.}
\loigiai{
\begin{itemize}
    \item Ý a) Sai: Kết quả đúng là $3$, không phải $2$. Khúc thứ nhất dài hơn khúc thứ hai số gang tay là: $9 - 6 = 3$ (gang tay).
\textbf{Đáp số:} $3$.
    \item Ý b) Sai: Kết quả đúng là $4$, không phải $5$. Số bạn nam nhiều hơn số bạn nữ là: $7 - 3 = 4$ (bạn).
\textbf{Đáp số:} $4$.
    \item Ý c) Đúng: Vì bạn Nam có số ô tô bằng bạn Tuấn, nên bạn Nam cũng nhiều hơn bạn Minh $2$ chiếc ô tô đồ chơi.
\textbf{Đáp số:} $2$.
    \item Ý d) Đúng: Túi thứ hai nhiều hơn túi thứ nhất số cái kẹo là: $9 - 5 = 4$ (cái kẹo).\nĐể hai túi có số kẹo bằng nhau, ta lấy số kẹo chênh lệch chia đều: $2 + 2 = 4$, vậy cần chuyển $2$ cái kẹo.\nKiểm tra lại: Túi 1 có $5 + 2 = 7$, túi 2 có $9 - 2 = 7$ (bằng nhau).
\textbf{Đáp số:} $2$.
\end{itemize}
\begin{itemize}
    \item a) \textbf{Sai}.
    \item b) \textbf{Sai}.
    \item c) \textbf{Đúng}.
    \item d) \textbf{Đúng}.
\end{itemize}
}
\end{ex}

\textbf{Phần 3. Trả lời ngắn}
\setcounterpageref{ex}{0}

\begin{ex} %[1D1V03]
Bạn Nga có 10 bông hoa, bạn Mai có 8 bông hoa. Hỏi bạn Nga phải tặng bạn Mai bao nhiêu bông hoa để hai bạn có số hoa bằng nhau?
\shortans{1}
\loigiai{
\textbf{Đáp số:} $1$\par
Số hoa bạn Nga nhiều hơn bạn Mai là: $10 - 8 = 2$ (bông hoa).\nĐể số hoa bằng nhau, Nga cần tặng Mai: $2 : 2 = 1$ bông hoa.\nKiểm tra: Khi đó Nga có $10 - 1 = 9$ bông, Mai có $8 + 1 = 9$ bông (bằng nhau).
\textbf{Đáp số:} $1$.
}
\end{ex}

\begin{ex} %[1D1V03]
Đàn gà nhà bác An có 6 con gà con và 1 con gà mẹ. Hỏi số gà con nhiều hơn số gà mẹ bao nhiêu con?
\shortans{5}
\loigiai{
\textbf{Đáp số:} $5$\par
Số gà con nhiều hơn gà mẹ là: $6 - 1 = 5$ (con gà).
\textbf{Đáp số:} $5$.
}
\end{ex}

\begin{ex} %[1D1V03]
Lan có 7 cái bút màu, Mai có ít hơn Lan 4 cái bút màu. Hỏi cả Lan và Mai có bao nhiêu cái bút màu?
\shortans{10}
\loigiai{
\textbf{Đáp số:} $10$\par
Số bút màu của Mai là: $7 - 4 = 3$ (cái bút).\nTổng số bút màu của cả hai bạn là: $7 + 3 = 10$ (cái bút).
\textbf{Đáp số:} $10$.
}
\end{ex}

\begin{ex} %[1D1V03]
Lớp 1B có 5 bạn thích vẽ tranh và 8 bạn thích múa hát. Hỏi số bạn thích múa hát nhiều hơn số bạn thích vẽ tranh bao nhiêu bạn?
\shortans{3}
\loigiai{
\textbf{Đáp số:} $3$\par
Số bạn thích múa hát nhiều hơn số bạn thích vẽ tranh là: $8 - 5 = 3$ (bạn).
\textbf{Đáp số:} $3$.
}
\end{ex}

\begin{ex} %[1D1V03]
Một giá sách có hai ngăn: ngăn trên có 10 quyển truyện, ngăn dưới có 5 quyển truyện. Hỏi cần chuyển bao nhiêu quyển truyện từ ngăn trên xuống ngăn dưới để ngăn trên vẫn nhiều hơn ngăn dưới 1 quyển?
\shortans{2}
\loigiai{
\textbf{Đáp số:} $2$\par
Nếu chuyển 2 quyển từ ngăn trên xuống ngăn dưới:\n- Ngăn trên còn: $10 - 2 = 8$ (quyển).\n- Ngăn dưới có: $5 + 2 = 7$ (quyển).\nKhi đó ngăn trên hơn ngăn dưới: $8 - 7 = 1$ (quyển, đúng yêu cầu đề bài).\nVậy cần chuyển $2$ quyển truyện.
\textbf{Đáp số:} $2$.
}
\end{ex}

\begin{ex} %[1D1V03]
Bạn Hùng có 4 quả bóng bàn, bạn Dũng có 10 quả bóng bàn. Hỏi bạn Dũng nhiều hơn bạn Hùng bao nhiêu quả bóng bàn?
\shortans{6}
\loigiai{
\textbf{Đáp số:} $6$\par
Số quả bóng bàn bạn Dũng nhiều hơn bạn Hùng là: $10 - 4 = 6$ (quả bóng bàn).
\textbf{Đáp số:} $6$.
}
\end{ex}
